\section{Evolution-Enhanced Graph-Attentive Cooperative Scheduling}
\label{sec:method}

\subsection{Framework and Dispatch Interface}
\label{sec:method-framework}
\label{sec:dispatch-interface}

Battery inventory couples dispatch across hours. Community exchange
couples microgrid costs. Graph-Attentive Value
Decomposition Actor-Critic (GAVDAC) uses recurrent histories to condition
community-value estimation and an evolutionary algorithm (EA) to explore
complete dispatch policies. Its policy-gradient backbone follows Factored
Multi-Agent Centralised Policy Gradients (FACMAC)~\cite{facmac2021};
graph mixing builds on GraphMIX~\cite{graphmix2020}, and population
interaction follows evolutionary reinforcement learning (ERL)~\cite{erl2018}.

Figure~\ref{fig:method-framework} shows the two learning branches.
The actor, local critic, and mixer have parameters $\theta$, $\phi$, and
$\psi$, respectively. A population of $M$ complete shared actors supplies
additional trajectories to replay buffer $\mathcal D$. Gradient learning
updates $(\theta,\phi,\psi)$; selection, mutation, and actor injection
update the population.

\begin{figure}[htbp]
 \centering
 \begin{tikzpicture}[
 x=1cm,y=1cm,
 font=\footnotesize,
 block/.style={draw=black!65,rounded corners=2pt,align=center,
   text width=3.05cm,minimum height=1.15cm,inner sep=4pt},
 flow/.style={-{Stealth[length=1.7mm]},line width=.65pt},
 update/.style={flow,dashed},
 lab/.style={font=\scriptsize,fill=white,inner sep=2pt,align=center}
]
 \node[block,fill=blue!7,minimum width=3.45cm] (pop) at (1.9,4.25)
   {Policy population\\$M$ shared actors};
 \node[block,fill=green!10,minimum width=3.45cm] (actor) at (9.5,4.25)
   {RL actor\\Shared policy $\theta$};
 \node[block,fill=black!3,text width=10.65cm,minimum width=10.95cm] (env) at (5.7,2.55)
   {Microgrid scheduling environment\\
    Common observations, feasible actions, and team reward};
 \node[block,fill=orange!8,minimum width=3.45cm] (evo) at (1.9,.70)
   {Evolution\\Selection, copying\\Mutation};
 \node[block,fill=blue!5,minimum width=3.25cm] (replay) at (5.7,.70)
   {Shared replay $\mathcal{D}$\\All trajectories};
 \node[block,fill=green!7,minimum width=3.45cm] (learner) at (9.5,.70)
   {GAVDAC learner\\Critic and graph mixer};

 \draw[flow] (pop.south)--(env.north -| pop.south)
   node[midway,lab] {Deterministic rollouts};
 \draw[flow] (actor.south)--(env.north -| actor.south)
   node[midway,lab] {Exploratory rollout};
 \draw[flow] (env.south)--(replay.north)
   node[midway,lab] {All training trajectories};
 \draw[flow] (replay.east)--(learner.west);
 \draw[flow] (env.south -| evo.north)--(evo.north)
   node[midway,lab] {Common-scenario fitness};

 \draw[update] (evo.west)--(-.10,.70)--(-.10,4.25)--(pop.west);
 \node[lab,rotate=90] at (-.10,2.55) {Next population};
 \draw[update] (learner.east)--(11.55,.70)--(11.55,4.25)--(actor.east);
 \node[lab,rotate=90] at (11.55,2.55) {Policy gradient};
 \draw[update] (actor.north)--(9.5,5.15)--(1.9,5.15)--(pop.north);
 \node[lab] at (5.7,5.15) {Periodic actor injection; elite protected};
 \draw[update] (learner.south)--(9.5,-.25)--(10.65,-.25)--(learner.south -| 10.65,0);
 \node[lab] at (9.35,-.48) {Critic and target updates};
\end{tikzpicture}

 \caption{Population exploration and graph-attentive learning. Solid arrows
 carry policies, actions, or experience; dashed arrows denote updates.}
 \label{fig:method-framework}
\end{figure}

Each agent receives local measurements, equipment limits, community
broadcasts, and causal forecast summaries as defined in
Section~\ref{sec:decision}. Its generator and battery coordinates pass
through Eq.~\eqref{eq:action-mapping}. Scaling battery actions by available
charging/discharging power exposes the feasible flexibility as SOC changes.
The centered reward retains the discounted-cost ranking of policies on
a common operating sequence. Exchange clearing and grid balancing follow
Section~\ref{sec:formulation}.

\subsection{History-Conditioned Graph Value Decomposition}
\label{sec:graph-decomposition}

The shared actor receives $[o_{i,t},a_{i,t-1},e_i]$, where $e_i$ is a
one-hot agent identifier. An input layer with a rectified linear unit (ReLU)
and a gated recurrent unit (GRU) produce hidden state $h_{i,t}$.
A linear output with a hyperbolic tangent gives the two-component vector
$\nu_{i,t}\in[-1,1]^2$, affinely mapped to action bounds
$[0,1]\times[-1,1]$. Parameters are shared, while observations,
identifiers, and hidden states distinguish microgrid roles.

A local critic gives utility $q_{i,t}=q_\phi(o_{i,t},a_{i,t})$.
History features are $\zeta_{i,t}=\operatorname{ReLU}(W_hh_{i,t}+b_h)$.
Attention scores are
$u_{ij,t}=(W_Q\zeta_{i,t})^\top(W_K\zeta_{j,t})/\sqrt{d_\zeta}$.
Here $W_h,b_h,W_Q,W_K$ are trainable projections and $d_\zeta$ is the
query/key dimension. For binary mask $C$, normalized weights are
\begin{equation}
 A_{ij,t}=\frac{C_{ij}\exp(u_{ij,t})}
                    {\sum_{k=1}^{N}C_{ik}\exp(u_{ik,t})}.
 \label{eq:mixer-attention}
\end{equation}

The implemented mask is
$C=\left[\begin{smallmatrix}1&1&1\\1&1&0\\1&0&1\end{smallmatrix}\right]$.
It defines utility-aggregation connections independently of physical
exchange links. Attention varies within this fixed mask as histories evolve.

For utility vector $q_t=(q_{1,t},\ldots,q_{N,t})^\top$, the mixed team
value is
\begin{equation}
 \begin{aligned}
 f_{i,t}&=g_\psi\big((A_tq_t)_i;s_t^{\mathrm{c}}\big),\\
 Q_{\mathrm{tot},t}
 &=\left(\frac1N\sum_i f_{i,t}\right)
          w_{\mathrm{out}}(s_t^{\mathrm{c}})+V(s_t^{\mathrm{c}}).
 \end{aligned}
 \label{eq:graph-value}
\end{equation}

The graph transformation $g_\psi$ applies two state-conditioned affine
layers with leaky ReLU and exponential linear unit activations.
The vector $f_{i,t}$ is the resulting node feature,
$w_{\mathrm{out}}$ weights its mean readout, and $V$ is a scalar bias.
These functions use the layer-normalized joint observation
$s_t^{\mathrm{c}}$. Hypernetworks generate nonnegative weights through
scaled sigmoids, following QMIX~\cite{qmix2018}. For fixed state and
recurrent features, the team value is nondecreasing in local utilities.
Actor histories thus condition both dispatch and community-value mixing.

\subsection{Joint-Value Policy Learning}
\label{sec:policy-learning}
\label{sec:facmac-learning}

Replay retains daily sequences to reconstruct recurrent histories.
Let $\tau_t$ denote joint replayed histories,
$\boldsymbol\mu_\theta$ the joint actor, and $\boldsymbol h_t^\theta$
its hidden states. Superscript $-$ denotes target networks.
For true-terminal indicator $d_t$, the bootstrap target is
\begin{equation}
 y_t=r_t+\gamma(1-d_t)Q_{\mathrm{tot}}^-
 \big(s_{t+1}^{\mathrm{c}},
      \boldsymbol\mu_{\theta^-}(\tau_{t+1}),\boldsymbol h_{t+1}^{-}\big).
 \label{eq:td-target}
\end{equation}

The critic and mixer minimize the mean squared target error over valid
transitions, with actor features detached. After this update, their
parameters are fixed while the actor minimizes
\begin{equation}
 \mathcal L_\pi(\theta)=
 -\left\langle Q_{\mathrm{tot}}
 (s_t^{\mathrm{c}},\boldsymbol\mu_\theta(\tau_t),\boldsymbol h_t^\theta)
 \right\rangle
 +\lambda_a\left\langle\|\nu_{i,t}\|_2^2/2\right\rangle.
 \label{eq:actor-loss}
\end{equation}

Angle brackets denote sample means, including agents for the action
regularizer with coefficient $\lambda_a$. Gradients pass through actions
and recurrent attention features. Target networks use soft parameter
updates after each learner update; settings appear in
Section~\ref{sec:experimental-setup}.

Previous actions in replay reconstruct both current and target histories.
Daily execution resets hidden states and previous actions, while SOC
continues across days. The final bootstrap extends the preceding history,
giving a blockwise approximation at the day boundary.

\subsection{Common-Scenario Evolution and Learning Interaction}
\label{sec:evolution}
\label{sec:hybrid-interaction}
\label{sec:training-procedure}

At each EA cycle, all $M$ candidates use deterministic actions on the same
training day with the same initial SOC vector. Candidate $k$ receives
fitness $F_k=\sum_{t=0}^{T-1}r_t^{(k)}$, its undiscounted daily team return.
Matching exogenous inputs and initial energy makes within-cycle fitness
differences attributable to policy behavior.

Rank-based tournaments select policies while protecting the elite.
The configured crossover copies one paired complete actor into the other,
with a randomly chosen direction. Mutation perturbs selected matrix
weights using relative Gaussian noise, with occasional large perturbations
or resets; bias vectors remain unchanged. One evaluated generation updates
the complete shared actor per cycle. Operator settings are given in
Section~\ref{sec:experimental-setup}.

All population trajectories enter $\mathcal D$. The gradient actor then
collects a trajectory from the same scenario using ordinary exploration;
its final SOC and date advance training. After replay warmup, learner
updates proceed in critic/mixer, actor, and target order. The updated actor
is injected into a population slot chosen from pre-evolution fitness while
protecting the retained elite. Population experience therefore reaches the
learner through replay, and gradient-trained behavior reaches the population
through injection.

Every population rollout counts toward the interaction budget. Validation
trajectories are excluded from replay, and the actor and population elite
remain separate checkpoint candidates. Online scheduling uses only the
selected deterministic actor, recurrent states, and prescribed
measurements/broadcasts, followed by feasible-power mapping and exchange
clearing. The critic, mixer, and evolution operate during training.
